\chapter{Conclusions and Future Work}
\label{chap:conclusions}

% =============================================================
\section{Conclusion}
\label{sec:conclusion}

The Vital Guardian project set out to address a critical and underserved problem
in modern healthcare: the absence of an intelligent, privacy-preserving monitoring
system for comatose and high-dependency ICU patients. Through a structured design
and development process, the project successfully produced a fully operational,
multi-modal monitoring prototype that demonstrates the viability of edge-based
AI in a clinical context.

\subsection{Visual Guardian}

The Visual Guardian module established a robust three-tier vision pipeline. The
YOLOv11n object detector, optimised via Intel OpenVINO, provides reliable real-time
patient localisation with minimal CPU overhead. The MoViNet-A2 spatiotemporal
classifiers operate on top of this localised region of interest, achieving a
video-level fall detection F1 score of 0.828 and a seizure detection recall of
88.0\% on the held-out evaluation sets. The MediaPipe Kinematic Safety Net
provides a rule-based secondary verification layer, suppressing false positives
generated by normal patient repositioning and by cases where the MoViNet model
fires on a completely motionless patient.

\subsection{Auditory Watchdog}

The Auditory Watchdog module delivers fully offline, real-time audio intelligence
without any dependency on external services. The dual-gated architecture, combining
Silero VAD and YAMNet, ensures that the system actively listens for non-verbal
distress sounds at all times while simultaneously protecting patient and visitor
privacy by automatically suppressing speech transcription during extended
conversations. The YAMNet Severity Taxonomy maps detected sounds to a
four-tier clinical severity scale, enabling the Cognitive Core to reason about
audio events with structured, quantified context rather than raw audio data.

\subsection{Cognitive Core}

The Cognitive Core demonstrated the most significant impact on overall system
accuracy. The ablation study confirms that introducing the two-tier Gemini
reasoning pipeline reduced the integrated system's false alarm rate from 6.0\%
(vision and audio alone) to 1.8\% in the final configuration, while simultaneously
improving overall accuracy to 97.2\%. The Tier-2 verification step acts as an
intelligent filter, and the Tier-3 enrichment step produces natural language
narratives that nursing staff can act upon immediately without requiring
interpretation of raw sensor data.

\subsection{Summary}

Collectively, the three modules realise the core project objective: a system that
reliably detects genuine patient emergencies, suppresses nuisance alarms, and
delivers contextually rich, human-readable alerts to clinical staff, all while
ensuring that no raw video or audio data leaves the edge device. The end-to-end
alert latency of 1.31 seconds satisfies the defined performance requirement of
under 2 seconds, confirming the system's suitability for time-critical ICU
deployment.

% =============================================================
\section{Future Work}
\label{sec:future_work}

The current prototype establishes a strong foundation. The following areas have
been identified as the most impactful directions for future development.

\subsection{Fully Localised Language Model}

The current Cognitive Core relies on Gemini language models for reasoning and
narrative generation. While only text-based event metadata is transmitted, a
future iteration should migrate this capability to a fully on-device small
language model, achieving a completely air-gapped deployment with no dependency
on any external network connectivity. Candidate architectures such as Llama-3 8B
and Phi-3 Mini are strong starting points given their demonstrated reasoning
capability at a fraction of the parameter count of full-scale models.

\subsection{Model Optimisation and Edge Acceleration}

The MoViNet-A2 models currently run at floating-point precision. Applying INT8
post-training quantisation and structured channel pruning via the Intel OpenVINO
toolkit is expected to reduce inference latency by 30 to 50 percent, enabling
deployment on lower-power embedded hardware such as the Intel NUC or dedicated
AI accelerator boards. This would significantly reduce the cost barrier for
hospital-wide deployment across multiple rooms.

\subsection{Formal Clinical Validation}

The current evaluation was conducted on simulation-based and publicly available
video datasets. A critical next step is formal validation in a controlled clinical
environment with the involvement of trained ICU nursing staff. Such a study would
benchmark alert accuracy and response time against real-world patient scenarios,
identify edge cases not covered by the current training data, and provide the
evidence base required for any regulatory submission.

\subsection{Hospital Information System Integration}

The current system delivers alerts exclusively to a local web dashboard.
Developing secure, standards-compliant interfaces (such as HL7 FHIR) to push
verified alerts and patient event logs directly into existing Hospital Information
Systems and Electronic Health Record platforms would allow Vital Guardian to
function as a native component of the hospital's clinical workflow rather than a
standalone tool.

\subsection{Predictive Analytics and Actigraphy}

The current system operates in a reactive mode, raising alerts only after a
critical event has been detected. The MediaPipe Sleep Restlessness sub-system
already accumulates kinetic energy data across the torso keypoints, forming the
basis for a longer-term actigraphy model. A future iteration could analyse trends
in patient movement cadence, restlessness frequency, and distress audio patterns
over multi-hour windows to generate proactive risk scores, alerting staff to
deteriorating patient conditions before a critical event occurs.
